\documentclass{beamer}

\usetheme{Madrid}
\title{Prompt Engineering}
\author{Mehmet Ali Akyol, PhD}
\institute{Ankara Medipol University}
\date{\today}

\begin{document}

\begin{frame}
    \titlepage
    \begin{center}
        \vspace{1cm}
        \textbf{Week 2: How LLMs Work (A Non-Technical Primer)}
    \end{center}
\end{frame}

\section{Overview}

\begin{frame}
    \frametitle{Weekly Learning Objectives}
    \begin{itemize}
        \item Explain an LLM as a next-word prediction engine using simple analogies.
    \item Define ``tokenization'' and its role in how LLMs process text.
        \item Describe the basic idea behind the Transformer architecture (self-attention).
        \item Distinguish between the capabilities and inherent limitations of LLMs.
    \end{itemize}
\end{frame}

\begin{frame}
    \frametitle{Key Concepts}
    \begin{itemize}
        \item Large Language Model (LLM)
        \item Transformer Architecture \& Self-Attention
        \item Tokenization \& Word Embeddings

        \item Capabilities \& Limitations
        \item Hallucination
    \end{itemize}
\end{frame}

\section{Core Concepts of LLMs}

\begin{frame}
    \frametitle{The Core Idea: A Sophisticated Autocomplete}
    \begin{itemize}
        \item Think of how your phone suggests the next word as you type—LLMs work on the same principle, just vastly more sophisticated.
    \item Instead of predicting just words, LLMs predict ``tokens'' (we'll explain this soon).
        \item They're trained by reading trillions of words from books, websites, code, scientific papers, and more.
        \item Through this reading, the model discovers patterns: grammar rules, common facts, writing styles, even cultural context.
    \item What's remarkable: complex behaviors like answering questions or writing essays emerge from this simple ``predict the next piece'' mechanism.
        \item \textbf{Key Insight:} No one explicitly programs these rules—the AI learns them by finding patterns in massive amounts of text.
    \end{itemize}
\end{frame}

\begin{frame}
    \frametitle{The Scale of Modern LLMs: How Big Are They?}
    \begin{itemize}
    \item \textbf{Parameters:} Think of these as the ``knobs'' the AI adjusts during learning. Modern LLMs have billions or even trillions of these adjustable values.
        \item \textbf{To put this in perspective:}
        \begin{itemize}
            \item GPT-3: 175 billion parameters (like 175 billion adjustable settings)
            \item GPT-4: Over 1 trillion parameters
            \item LLaMA 2 70B: 70 billion parameters
        \end{itemize}
        \item \textbf{Training Data:} These models read hundreds of billions of text pieces during training—imagine reading millions of books.
        \item \textbf{Cost:} Training cutting-edge models can cost millions of dollars in computing power and electricity.
        \item Generally, bigger models perform better, though there are limits to this trend.
    \end{itemize}
\end{frame}

\begin{frame}
    \frametitle{Small vs. Large Models: What's the Difference?}
    \begin{itemize}
        \item Think of it like comparing a high school student to a college professor—both can answer questions, but with different depth.
        \item \textbf{Small Models (fewer parameters):}
        \begin{itemize}
            \item Faster and cheaper to run
            \item Good for simple tasks (basic Q\&A, simple summaries)
            \item May miss nuances or make more mistakes
            \item Example response: ``Paris is the capital of France. It has many tourists.''
        \end{itemize}
        \item \textbf{Large Models (more parameters):}
        \begin{itemize}
            \item Better at complex reasoning and nuanced understanding
            \item Can handle specialized topics and longer contexts
            \item Example response: ``Paris, France's capital, is renowned for its art, culture, and history. The Eiffel Tower, Louvre Museum, and Notre-Dame Cathedral attract millions of visitors annually\ldots''
        \end{itemize}
    \end{itemize}
\end{frame}

\begin{frame}
    \frametitle{From Words to Numbers: Tokenization}
    \begin{itemize}
    \item Here's a crucial point: computers can't actually ``read'' text---they only understand numbers.
    \item \textbf{Tokenization} is how we bridge this gap: breaking text into small chunks called ``tokens.''
        \item A token might be:
        \begin{itemize}
            \item A complete word (like ``hello'')
            \item Part of a word (like ``un'' and ``believable'')
            \item Even a single character or punctuation mark
        \end{itemize}
    \item \textbf{Example:} The phrase ``Tokenization is important'' might be split into: [``Token'', ``ization'', `` is'', `` important''].
        \item Each token is then converted to a number the computer can process.
    \end{itemize}
\end{frame}

\begin{frame}
    \frametitle{Why Break Words Into Pieces?}
    \begin{itemize}
        \item \textbf{The Solution:} Modern LLMs use 32,000-100,000 ``sub-word'' pieces that can be combined flexibly.
        \item \textbf{Why This Works:}
        \begin{itemize}
            \item Can handle new or rare words by combining familiar pieces (like building with LEGO blocks)
            \item More efficient for the computer to process
            \item Works especially well for languages like Turkish or Finnish where words can be very long
        \end{itemize}
        \item Different AI models may split text differently—like different people cutting a cake into different slice sizes.
        \item This is why the same text might be slightly different lengths for different AI systems.
    \end{itemize}
\end{frame}

\begin{frame}
    \frametitle{From Tokens to Meaning: Embeddings}
    \begin{itemize}
        \item After splitting text into tokens, each token becomes a list of numbers called an \textbf{embedding}.
        \item Think of it like GPS coordinates: just as latitude and longitude locate a place on Earth, embeddings locate words in ``meaning space.''
        \item Each token gets represented by 768 to 12,288 numbers that capture its meaning and relationships to other words.
        \item \textbf{The Magic:} Words with similar meanings end up close together in this mathematical space.
        \item \textbf{Famous Example:} The relationship between "king" and "queen" is mathematically similar to the relationship between "man" and "woman."
        \item The AI learns these relationships from data—no human explicitly teaches it that kings and queens are related!
    \end{itemize}
\end{frame}
\begin{frame}
    \frametitle{Visualizing Embeddings: Words in Space}
    \begin{center}
        \textbf{Imagine a Map of Word Meanings}
    \end{center}
    \begin{itemize}
        \item If we could visualize this ``meaning space,'' related words would appear in neighborhoods together:
        \begin{itemize}
            \item Animal neighborhood: cat, dog, horse, elephant
            \item Color district: red, blue, green, yellow
            \item Action zone: run, walk, jump, sprint
        \end{itemize}
    \item \textbf{Context Matters:} Modern LLMs are smart---they understand that ``bank'' means something different in:
        \begin{itemize}
            \item ``I walked along the river bank'' (shoreline)
            \item ``I deposited money at the bank'' (financial institution)
        \end{itemize}
    \item The same word gets different embeddings based on context!
    \end{itemize}
\end{frame}

\begin{frame}
    \frametitle{The Brain of the LLM: The Transformer}
    \begin{itemize}
        \item In 2017, researchers published a breakthrough paper titled ``Attention Is All You Need'' that changed everything.
        \item The key innovation: \textbf{self-attention}---a way for AI to understand which words are important for understanding other words.
        \item \textbf{Think about this sentence:} ``The cat sat on the mat, and \textbf{it} was happy.''
        \item You instantly know ``it'' refers to ``the cat'', not ``the mat''. How? You're paying attention to the relationship between words.
        \item Self-attention lets the AI do the same thing: mathematically determine which words should influence the understanding of other words.
        \item This replaced older approaches that read text one word at a time, making modern AI both faster and smarter.
    \end{itemize}
\end{frame}

\begin{frame}
    \frametitle{How Self-Attention Works (Simplified)}
    \begin{itemize}
    \item Imagine reading a sentence where each word asks: ``Which other words should I pay attention to?''
        \item For each word, the system creates three pieces of information:
        \begin{itemize}
            \item \textbf{Query:} ``What am I looking for?''
            \item \textbf{Key:} ``What information do I contain?''
            \item \textbf{Value:} ``What information can I provide?''
        \end{itemize}
    \item The system compares all the Queries and Keys to calculate \textbf{attention scores}---essentially asking ``how relevant is each word to each other word?''
        \item These scores determine how much each word influences the understanding of every other word.
        \item \textbf{The Result:} Each word gets a rich, context-aware representation that incorporates information from the entire sentence.
    \end{itemize}
\end{frame}

\begin{frame}
    \frametitle{Multi-Head Attention: Many Perspectives at Once}
    \begin{itemize}
        \item Here's where it gets even more interesting: the AI doesn't use just one attention mechanism—it uses many in parallel.
    \item \textbf{Example:} GPT-3 uses 96 different ``attention heads'' in each layer, and has 96 layers! That's over 9,000 attention mechanisms working together.
        \item Think of it like having multiple experts examining the same text:
        \begin{itemize}
            \item One expert focuses on grammar (subject-verb agreement, sentence structure)
            \item Another on meaning (synonyms, related concepts)
            \item Another on long-distance connections (what ``it'' refers to paragraphs earlier)
        \end{itemize}
        \item All these perspectives get combined to create a rich, nuanced understanding of the text.
    \end{itemize}
\end{frame}

\begin{frame}
    \frametitle{The Transformer Architecture: Putting It All Together}
    \begin{itemize}
        \item \textbf{Step 1 - Input:} Text is converted to embeddings (those number lists we discussed)
        \item \textbf{Step 2 - Position Information:} The system adds information about word order (first word, second word, etc.)
        \item \textbf{Step 3 - Processing Layers:} The text flows through many layers (12 to 96+ in modern models)
        \begin{itemize}
            \item Each layer applies multi-head attention (understanding word relationships)
            \item Then applies additional processing to refine understanding
        \end{itemize}
        \item \textbf{Step 4 - Output:} Finally, the system produces probabilities for what token should come next
        \item Think of it as multiple stages of increasingly sophisticated reading comprehension.
    \end{itemize}
\end{frame}

\section{Training and Fine-Tuning}

\begin{frame}
    \frametitle{The Training Process: How Do They Learn?}
    \begin{itemize}
        \item \textbf{Pre-training:} The AI reads massive amounts of text—books, websites, articles, code, and more.
        \item \textbf{The Learning Game:} For each piece of text, try to predict the next word. Get it wrong? Adjust those billions of parameters slightly. Repeat trillions of times.
        \item It's like learning a language by reading millions of books and constantly testing yourself on what word comes next.
        \item \textbf{Time \& Resources:} This process takes weeks to months, running on thousands of powerful computers working together.
        \item \textbf{The Result:} A "foundation model"—an AI with broad knowledge about language, facts, reasoning patterns, and more.
    \end{itemize}
\end{frame}

\begin{frame}
    \frametitle{Fine-Tuning: Making Them Useful and Safe}
    \begin{itemize}
        \item After basic training, the AI isn't ready for public use—it needs refinement.
        \item \textbf{Fine-Tuning:} Additional training on carefully selected examples to specialize the model's behavior.
        \item \textbf{Instruction Tuning:} Teaching the AI to follow instructions by showing it many examples of questions and good answers.
        \item \textbf{Learning from Human Feedback:}
        \begin{itemize}
            \item Humans review and rank different AI responses
            \item The AI learns to produce responses that humans find helpful, harmless, and honest
            \item This helps align the AI with human values and expectations
        \end{itemize}
        \item Systems like ChatGPT, Claude, and Gemini all go through extensive fine-tuning before public release.
    \end{itemize}
\end{frame}

\section{Capabilities, Limitations, and Activities}

\begin{frame}
    \frametitle{Capabilities and Limitations}
    \begin{block}{What LLMs Are Remarkably Good At:}
        LLMs are remarkably flexible—they can write, summarize, translate, analyze, code, and more without being explicitly programmed for each task. This general-purpose adaptability is revolutionary.
    \end{block}
    
    \begin{block}{Key Limitations to Keep in Mind:}
        \begin{itemize}
            \item Hallucination: can produce plausible but incorrect or fabricated facts.
            \item Biases: reflect stereotypes and biases present in training data.
            \item No-true understanding: statistical pattern learners, not conscious reasoners.
            \item Outdated/incomplete knowledge: training cutoff and dataset gaps.
            \item Context-window limits: cannot attend to arbitrarily long documents.
            \item Sensitivity to phrasing: small prompt changes produce different outputs.
            \item Privacy \& safety risks: may leak training artifacts or be misused.
        \end{itemize}
    \end{block}
\end{frame}

\begin{frame}
    \frametitle{Emergent Abilities: Unexpected Superpowers}
    \begin{itemize}
        \item Something fascinating happens as models get bigger: new abilities suddenly appear that smaller models didn't have.
        \item \textbf{Examples of emergent abilities:}
        \begin{itemize}
            \item Learning from just a few examples (like humans do)
            \item Explaining step-by-step reasoning
            \item Following complex, multi-part instructions
            \item Solving problems they've never seen before
        \end{itemize}
        \item Here's what's remarkable: nobody explicitly programmed these abilities—they spontaneously emerged when the models reached a certain size.
        \item \textbf{Open Question:} Are these truly new abilities, or are we just measuring things differently?
        \item Understanding this helps us anticipate what even larger future models might be capable of.
    \end{itemize}
\end{frame}

\begin{frame}
    \frametitle{In-Class Activities}
    \begin{block}{Activity 1: Breaking Down Sentences Into Tokens}
        Let's experiment with tokenization! We'll take several sentences and first predict how an AI might break them into tokens. Then we'll use an online tokenizer tool (platform.openai.com/tokenizer) to see the actual breakdown. How close were our predictions? What surprises us?
    \end{block}
    \begin{block}{Activity 2: Spot the Hallucination}
        Time to be AI detectives! We'll examine AI-generated paragraphs that contain subtle factual errors. Using fact-checking and search engines, we'll work together to identify which information is hallucinated and discuss why the AI might have made these mistakes.
    \end{block}
\end{frame}

% --- Activity 2: Spot the Hallucination ---
% AI & Technology
\begin{frame}{Activity 2: Spot the Hallucination}
    	\textbf{Domain:} [AI] AI \& Technology
    \vspace{0.3cm}
    	\textit{Hallucination:}
    \begin{quote}
        ``OpenAI released GPT-5 in 2018, and it quickly became the first model capable of passing the Turing test and operating fully without any training data''.
    \end{quote}
    \vspace{0.3cm}
    	\textbf{Can you spot the error(s)?}
\end{frame}

\begin{frame}{Activity 2: AI \& Technology — Reveal}
    	\textbf{What’s Wrong:}
    \begin{itemize}
        \item GPT-5 wasn’t released in 2018 (GPT-2 came in 2019).
        \item No AI model has definitively passed the Turing test.
        \item All models require training data.
    \end{itemize}
    	\textbf{Correct Fact:}
    \begin{itemize}
        \item GPT-2 was released in 2019, GPT-3 in 2020, GPT-4 in 2023.
        \item None have conclusively passed the Turing test.
        \item All are trained on large datasets.
    \end{itemize}
\end{frame}

% Medical (1)
\begin{frame}{Activity 2: Spot the Hallucination}
    	\textbf{Domain:} [MED] Medical
    \vspace{0.3cm}
    	\textit{Hallucination:}
    \begin{quote}
        ``The hormone insulin was first discovered in the 1970s by Edward Jenner, who used it to successfully treat smallpox patients''.
    \end{quote}
    \vspace{0.3cm}
    	\textbf{Can you spot the error(s)?}
\end{frame}

\begin{frame}{Activity 2: Medical — Reveal}
    	\textbf{What’s Wrong:}
    \begin{itemize}
        \item Insulin was discovered in 1921 by Banting and Best.
        \item Jenner lived much earlier and developed the smallpox vaccine.
    \end{itemize}
    	\textbf{Correct Fact:}
    \begin{itemize}
        \item Insulin was discovered in 1921 and is used to treat diabetes.
        \item Jenner developed the smallpox vaccine in the late 18th century.
    \end{itemize}
\end{frame}

% Medical (2)
\begin{frame}{Activity 2: Spot the Hallucination}
    	\textbf{Domain:} [MED] Medical
    \vspace{0.3cm}
    	\textit{Hallucination:}
    \begin{quote}
        ``The human heart pumps blue blood to the lungs, where it turns red after mixing with carbon dioxide, which provides oxygen to body tissues''.
    \end{quote}
    \vspace{0.3cm}
    	\textbf{Can you spot the error(s)?}
\end{frame}

\begin{frame}{Activity 2: Medical — Reveal}
    	\textbf{What’s Wrong:}
    \begin{itemize}
        \item Blood is never blue; it just appears that way under the skin.
        \item Blood releases CO$_2$ in the lungs and binds oxygen — it doesn’t mix with CO$_2$ to gain oxygen.
    \end{itemize}
    	\textbf{Correct Fact:}
    \begin{itemize}
        \item Deoxygenated blood is dark red and becomes bright red when oxygen binds to hemoglobin in the lungs.
    \end{itemize}
\end{frame}

% Space & Astronomy
\begin{frame}{Activity 2: Spot the Hallucination}
    	\textbf{Domain:} [SPACE] Space \& Astronomy
    \vspace{0.3cm}
    	\textit{Hallucination:}
    \begin{quote}
        ``Jupiter is the smallest planet in our solar system and completes one orbit around the Sun every 12 days, which makes its years much shorter than Earth’s''.
    \end{quote}
    \vspace{0.3cm}
    	\textbf{Can you spot the error(s)?}
\end{frame}

\begin{frame}{Activity 2: Space \& Astronomy — Reveal}
    	\textbf{What’s Wrong:}
    \begin{itemize}
        \item Jupiter is the largest planet, not the smallest.
        \item Its orbital period is about 12 Earth years, not 12 days.
    \end{itemize}
    	\textbf{Correct Fact:}
    \begin{itemize}
        \item Jupiter is the largest planet in the solar system and orbits the Sun roughly every 12 Earth years.
    \end{itemize}
\end{frame}

% History
\begin{frame}{Activity 2: Spot the Hallucination}
    	\textbf{Domain:} [HIST] History
    \vspace{0.3cm}
    	\textit{Hallucination:}
    \begin{quote}
        ``The Berlin Wall was built in 1945 immediately after World War II to separate democratic East Germany from communist West Germany''.
    \end{quote}
    \vspace{0.3cm}
    	\textbf{Can you spot the error(s)?}
\end{frame}

\begin{frame}{Activity 2: History — Reveal}
    	\textbf{What’s Wrong:}
    \begin{itemize}
        \item The Berlin Wall was built in 1961, not 1945.
        \item It separated communist East Berlin (and East Germany) from democratic West Berlin.
    \end{itemize}
    	\textbf{Correct Fact:}
    \begin{itemize}
        \item The wall was constructed in 1961 by East Germany to stop citizens from fleeing to West Berlin.
    \end{itemize}
\end{frame}

% Economy
\begin{frame}{Activity 2: Spot the Hallucination}
    	\textbf{Domain:} [ECON] Economy
    \vspace{0.3cm}
    	\textit{Hallucination:}
    \begin{quote}
        ``The Great Depression began in 1969 after the collapse of the European Union’s single currency, causing global GDP to shrink by 70\% within a single year''.
    \end{quote}
    \vspace{0.3cm}
    	\textbf{Can you spot the error(s)?}
\end{frame}

\begin{frame}{Activity 2: Economy — Reveal}
    	\textbf{What’s Wrong:}
    \begin{itemize}
        \item The Great Depression began in 1929.
        \item The EU and the euro did not exist then.
        \item Global GDP didn’t shrink by 70\% in one year — the U.S. GDP dropped ~30\% over several years.
    \end{itemize}
    	\textbf{Correct Fact:}
    \begin{itemize}
        \item The Great Depression started in 1929 after the U.S. stock market crash.
        \item The EU and euro were established decades later.
    \end{itemize}
\end{frame}

% End of Activity 2

\begin{frame}
    \frametitle{Discussion Questions}
    \begin{itemize}
        \item If an LLM is just predicting the next word, why does it feel like it understands our questions?
        \item Now that we know LLMs can hallucinate, how should this change the way we use them for academic research?
        \item What dangers might arise if people aren't aware of the biases embedded in LLMs?
        \item Does the concept of "emergent abilities" challenge our understanding of what intelligence really is?
        \item If LLMs can only "read" limited amounts of text at once, what creative strategies could we use to analyze very long documents?
        \item Should we trust AI-generated content differently than human-written content? Why or why not?
    \end{itemize}
\end{frame}


\begin{frame}
    \frametitle{Key Takeaways}
    \begin{itemize}
        \item LLMs are sophisticated pattern-matching systems that learn from enormous amounts of text—think of them as "autocomplete on steroids."
        \item Tokens are the building blocks of text (like LEGO bricks), and understanding tokenization helps explain why AI sometimes behaves unexpectedly.
        \item The Transformer architecture's self-attention mechanism is what enables them to understand context and relationships between words.
        \item Bigger models generally perform better, but there are trade-offs in cost, speed, and complexity.
        \item LLMs are powerful tools, but they have real limitations—hallucinations, biases, and knowledge gaps—that we must be aware of.
        \item Context windows limit how much text an LLM can "see" at once—this matters for long documents.
        \item Understanding how LLMs work helps us write better prompts and use AI more effectively in our own fields.
    \end{itemize}
\end{frame}

\begin{frame}
    \frametitle{Reflection for Next Week}
    \begin{block}{Think About This Before Our Next Class}
        \textbf{Reflection Prompt:} "If LLMs are just predicting the next word based on patterns, does that mean they're truly 'intelligent' or just very sophisticated calculators?"
        
        \vspace{0.3cm}
        \textbf{Consider:}
        \begin{itemize}
            \item What does 'understanding' really mean?
            \item Can pattern recognition alone produce intelligent behavior?
            \item How is this different from (or similar to) how humans learn and think?
            \item What are the philosophical implications for AI in society?
        \end{itemize}
        \vspace{0.3cm}
        \textbf{Next Class:} We'll share our thoughts and perspectives.
    \end{block}
\end{frame}

% Thank You Slide
\begin{frame}{Questions and Discussion}
    \begin{center}
    {\Huge Thank You!}
    \vspace{1cm}

    	\textbf{Dr.\ Mehmet Ali Akyol}\\
    	\texttt{mehmetali.akyol@gmail.com}
    \vspace{1cm}
    
    {\large Questions?}
    \end{center}
\end{frame}
\end{document}